\documentclass[a4paper,12pt]{article}
\usepackage[utf8]{inputenc}
\usepackage{ctex}
\usepackage{geometry}
\usepackage{listings}
\usepackage{xcolor}
\usepackage{amsmath}
\usepackage{amssymb} % 提供 \checkmark
\usepackage{pifont}  % 提供 \ding
\usepackage{hyperref}
\usepackage{enumitem}
\usepackage{fancyhdr}
\usepackage{titlesec}
\usepackage{tcolorbox}
\usepackage{parskip} % 段落间距，不缩进
\usepackage{tikz}
\usetikzlibrary{arrows, positioning, shapes}

% 页面边距设置
\geometry{left=2.5cm,right=2.5cm,top=2.5cm,bottom=2.5cm}

% 页眉页脚设置
\pagestyle{fancy}
\fancyhf{}
\fancyhead[L]{C语言与计算思维复习概要}
\fancyhead[R]{\leftmark}
\fancyfoot[C]{\thepage}
\renewcommand{\headrulewidth}{0.4pt}
\renewcommand{\footrulewidth}{0pt}

% 章节标题格式
\titleformat{\section}
  {\normalfont\Large\bfseries\color{blue!70!black}}{\thesection}{1em}{}[{\titlerule[0.8pt]}]
\titleformat{\subsection}
  {\normalfont\large\bfseries\color{blue!60!black}}{\thesubsection}{1em}{}

% 代码高亮设置
\definecolor{codegreen}{rgb}{0,0.6,0}
\definecolor{codegray}{rgb}{0.5,0.5,0.5}
\definecolor{codepurple}{rgb}{0.58,0,0.82}
\definecolor{backcolour}{rgb}{0.96,0.96,0.96}

\lstdefinestyle{mystyle}{
    backgroundcolor=\color{backcolour},   
    commentstyle=\color{codegreen},
    keywordstyle=\color{magenta}\bfseries,
    numberstyle=\tiny\color{codegray},
    stringstyle=\color{codepurple},
    basicstyle=\ttfamily\footnotesize,
    breakatwhitespace=false,         
    breaklines=true,                 
    captionpos=b,                    
    keepspaces=true,                 
    numbers=left,                    
    numbersep=5pt,                  
    showspaces=false,                
    showstringspaces=false,
    showtabs=false,                  
    tabsize=4,
    frame=single, % 代码框
    rulecolor=\color{gray!30},
    language=C
}
\lstset{style=mystyle}

% 超链接颜色设置
\hypersetup{
    colorlinks=true,
    linkcolor=blue!70!black,
    filecolor=magenta,      
    urlcolor=cyan,
    pdftitle={C语言复习攻略},
    pdfpagemode=FullScreen,
}

% 列表间距优化
\setlist[itemize]{itemsep=2pt, topsep=2pt, parsep=0pt}
\setlist[enumerate]{itemsep=2pt, topsep=2pt, parsep=0pt}

% 定义重点提示框
\newtcolorbox{tipbox}[1][]{colback=blue!5!white,colframe=blue!75!black,title=核心考点/易错点,fonttitle=\bfseries,#1}

\title{\textbf{\Huge C语言与计算思维\\期末复习全攻略}}
\author{\Large 摆渡 (Powered by GitHub Copilot)}
\date{\today}

\begin{document}

\maketitle
\thispagestyle{empty} % 首页不显示页码

\begin{center}
    \begin{tcolorbox}[colback=green!5!white,colframe=green!50!black,title=使用建议,fonttitle=\bfseries]
        本复习攻略旨在帮助大家高效梳理《C语言与计算思维》课程的核心考点。
        
        \textbf{建议用法}：请\textbf{快速浏览}全文，对于已经掌握的知识点一带而过，重点关注\textbf{自己不熟悉}或\textbf{容易混淆}的部分（如指针、文件操作、易错点辨析等）。
        
        \vspace{0.5em}
        \centering\textbf{\large 祝大家考试顺利，高分通过！}
    \end{tcolorbox}
\end{center}

\newpage
\tableofcontents
\newpage

\section{计算机系统与原理}

\subsection{核心考点}
\begin{itemize}
    \item \textbf{图灵测试 (Turing Test)}：测试机器是否具有人类智能的标准。如果测试者在不接触被测试者的情况下，无法区分人和机器，则认为机器通过测试。
    \item \textbf{计算机系统组成}：
    \begin{itemize}
        \item \textbf{硬件 (Hardware)}：CPU、存储器、输入/输出设备（冯·诺依曼结构核心）。
        \item \textbf{软件 (Software)}：系统软件（OS、驱动）、应用软件（Office、游戏）。
    \end{itemize}
    \item \textbf{冯·诺依曼体系结构}：
    \begin{itemize}
        \item \textbf{核心思想}：\textbf{存储程序} (Stored Program)。程序和数据都以二进制形式存储在存储器中。
        \item \textbf{五大部件}：运算器、控制器、存储器、输入设备、输出设备。（\hyperref[sec:von_neumann]{点击查看详细原理与图解}）
    \end{itemize}
    \item \textbf{数据表示}：
    \begin{itemize}
        \item \textbf{二进制}：计算机内部使用二进制（0和1）。
        \item \textbf{原码、反码与补码}：
        \begin{itemize}
            \item \textbf{原码}：最高位为符号位（0正1负），其余位表示数值大小。
            \item \textbf{反码}：正数的反码同原码；负数的反码为原码符号位不变，其余位按位取反。
            \item \textbf{补码}：正数的补码同原码；负数的补码为反码加1。
            \item \textbf{存储规则}：\textbf{负数在计算机中以补码形式存储}。
            \item \textbf{示例}（以8位二进制数 $-5$ 为例）：
            \begin{itemize}
                \item 原码：\texttt{1000 0101} （最高位1表示负，数值5为101）
                \item 反码：\texttt{1111 1010} （符号位不变，数值位取反）
                \item 补码：\texttt{1111 1011} （反码 + 1）
            \end{itemize}
        \end{itemize}
    \end{itemize}
\end{itemize}

\subsection{简答题预测}
\begin{description}
    \item[Q: 什么是“存储程序”原理？] \hfill \\
    A: 将程序（指令序列）和数据一起存储在计算机的存储器中，计算机在运行时自动地从存储器中取出指令并执行，而不需要人工干预。
    \item[Q: 简述硬件与软件的关系。] \hfill \\
    A: 硬件是计算机的躯壳（物理实体），软件是计算机的灵魂（思想）。硬件提供物质基础，软件指挥硬件工作。两者相辅相成，缺一不可。
\end{description}

\section{算法基础}

\subsection{核心考点}
\begin{itemize}
    \item \textbf{定义}：解决特定问题的一系列有限步骤。
    \item \textbf{五大特性}：
    \begin{enumerate}
        \item \textbf{有穷性}：步骤有限。
        \item \textbf{确定性}：指令清晰无歧义。
        \item \textbf{输入}：0个或多个。
        \item \textbf{输出}：1个或多个。
        \item \textbf{可行性}：每一步都能通过基本运算实现。
    \end{enumerate}
\end{itemize}

\subsection{简答题预测}
\begin{description}
    \item[Q: 什么是算法？它有哪些基本特征？] \hfill \\
    A: 算法是解决特定问题的一系列明确、有限的指令。基本特征包括：有穷性、确定性、可行性、输入、输出。
\end{description}

\section{实验与期末试卷分析}

\subsection{实验课核心内容总结}
\begin{itemize}
    \item \textbf{基础运算与日期}：
    \begin{itemize}
        \item 模拟计算器：注意运算符优先级和除零错误。
        \item 日期计算：闰年判断 \texttt{(year\%4==0 \&\& year\%100!=0) || (year\%400==0)}。
    \end{itemize}
    \item \textbf{数组与结构体}：
    \begin{itemize}
        \item 有序数组 vs 无序数组：插入、删除、查找的时间复杂度对比。
        \item 结构体数组：如客户信息管理、歌唱比赛评分（去掉最高最低分求平均）。
    \end{itemize}
    \item \textbf{递归应用}：
    \begin{itemize}
        \item 汉诺塔：经典的递归分治问题。
        \item 预测赢家：博弈论中的递归求解。
    \end{itemize}
    \item \textbf{大数运算}：
    \begin{itemize}
        \item 使用字符数组或整型数组模拟超大整数的加法，注意进位处理。
    \end{itemize}
    \item \textbf{栈的应用}：
    \begin{itemize}
        \item JSON字符串括号匹配：利用栈的后进先出特性，遇到左括号入栈，遇到右括号出栈匹配。
    \end{itemize}
\end{itemize}

\subsection{期末试卷高频考点解析}
\begin{itemize}
    \item \textbf{选择题陷阱与对比辨析}：
    \begin{itemize}
        \item \textbf{字符串长度}：\texttt{strlen} 遇到 \texttt{\textbackslash 0} 停止，转义字符（如 \texttt{\textbackslash 101}）算一个字符。
        \begin{itemize}
            \item \textbf{对比}：\texttt{sizeof} 统计内存空间大小（含 \texttt{\textbackslash 0}）；\texttt{strlen} 统计有效字符长度（不含 \texttt{\textbackslash 0}）。
            \item \textbf{例子}：
\begin{lstlisting}
char s[] = "hi";
// sizeof(s) -> 3
// strlen(s) -> 2
\end{lstlisting}
        \end{itemize}
        \item \textbf{逻辑短路}：\texttt{a || b++}，若 \texttt{a} 为真，\texttt{b++} 不执行。
        \begin{itemize}
            \item \textbf{对比}：\texttt{\&\&} 左假则右不执行；\texttt{||} 左真则右不执行。位运算 \texttt{\&} 和 \texttt{|} 无此特性。
            \item \textbf{例子}：
\begin{lstlisting}
int a = 1, b = 0;
if (a || b++) { ... }
// 执行后 b 仍为 0
\end{lstlisting}
            \item \textbf{真题陷阱}：
\begin{lstlisting}
int a = 1, b = 2, m = 0, n = 0, k;
k = (n = b > a) || (m = a < b);
// 解析：
// 1. b > a (2 > 1) 为真(1)，赋给 n，n 变为 1。
// 2. 左侧表达式 (n=...) 结果为 1 (真)。
// 3. || 发生短路，右侧 (m=a<b) 不执行！m 保持 0。
// 4. k 得到 || 的结果 1。
\end{lstlisting}
        \end{itemize}
        \item \textbf{指针类型}：区分 \texttt{int *p[10]} (指针数组) 和 \texttt{int (*p)[10]} (数组指针)。
        \begin{itemize}
            \item \textbf{对比}：前者是\textbf{数组}（存指针）；后者是\textbf{指针}（指数组）。
            \item \textbf{例子}：
\begin{lstlisting}
// 指针数组
char *strs[3] = {"A", "B", "C"};
// 数组指针
int arr[10];
int (*p)[10] = &arr;
\end{lstlisting}
        \end{itemize}
        \item \textbf{作用域}：局部变量屏蔽同名全局变量。
        \begin{itemize}
            \item \textbf{辨析}：同名时，谁近用谁（局部优先）。
            \item \textbf{例子}：
\begin{lstlisting}
int x = 10;
void f() {
    int x = 5;
    printf("%d", x); // 输出 5
}
\end{lstlisting}
        \end{itemize}
    \end{itemize}
    \item \textbf{程序填空题常考模式}：
    \begin{itemize}
        \item 详见附录 \hyperref[sec:appendix_code]{程序填空题常考模式代码详解}。
    \end{itemize}
    \item \textbf{代码阅读题}：
    \begin{itemize}
        \item 能够正确跟踪指针的移动和数组元素的访问。
        \item 理解 \texttt{sizeof} 在不同情况下的值（数组名 vs 指针）。
        \begin{itemize}
            \item \textbf{对比}：数组名在定义域内 \texttt{sizeof} 为总字节数；传参后退化为指针，\texttt{sizeof} 为指针大小（4或8）。
            \item \textbf{例子}：
\begin{lstlisting}
int a[10]; // sizeof(a) -> 40
void f(int a[]) {
    // sizeof(a) -> 4/8 (指针)
}
\end{lstlisting}
        \end{itemize}
    \end{itemize}
\end{itemize}

\section{C语言基础语法}

\subsection{数据类型与运算}
\label{sec:data_types}
\begin{itemize}
    \item \textbf{数据类型大小}：\texttt{sizeof} 是运算符。
    \begin{itemize}
        \item \texttt{char} (1), \texttt{short} (2), \texttt{int} (4), \texttt{long} (4或8), \texttt{double} (8)。
    \end{itemize}
    \item \textbf{类型转换}：
    \begin{itemize}
        \item \textbf{隐式}：\texttt{int + double} $\rightarrow$ \texttt{double}。
        \item \textbf{强制}：\texttt{(int)(3.14 + 0.9)} 结果是 4。
    \end{itemize}
    \item \textbf{自增自减}：
    \begin{itemize}
        \item \texttt{b = a++}：先赋值，后自增。
        \item \texttt{b = ++a}：先自增，后赋值。
        \item \textbf{真题陷阱}：\texttt{k=(n=b>a)||(m=a<b)}。注意逻辑短路，若前半部分为真，后半部分不执行。
    \end{itemize}
    \item \textbf{位运算}：
    \begin{itemize}
        \item \texttt{\&, |, \^{}, \~{}, <<, >>}。
        \item 技巧：\texttt{a << 1} (乘2), \texttt{a >> 1} (除2), \texttt{n \& 1} (判断奇偶)。
    \end{itemize}
    \item \textbf{运算符优先级}（由高到低）：
    \begin{enumerate}
        \item \textbf{初等}：\texttt{()} \texttt{[]} \texttt{->} \texttt{.}
        \item \textbf{单目}：\texttt{!} \texttt{\~} \texttt{++} \texttt{--} \texttt{sizeof} \texttt{(type)} \texttt{*} \texttt{\&} \hfill \textit{(右到左)}
        \item \textbf{算术}：\texttt{*} \texttt{/} \texttt{\%} $\rightarrow$ \texttt{+} \texttt{-}
        \item \textbf{移位}：\texttt{<<} \texttt{>>}
        \item \textbf{关系 (大小)}：\texttt{>} \texttt{<} \texttt{>=} \texttt{<=}
        \item \textbf{关系 (相等)}：\texttt{==} \texttt{!=} \hfill \textit{(优先级低于大小比较)}
        \item \textbf{位逻辑}：\texttt{\&} $\rightarrow$ \texttt{\^} $\rightarrow$ \texttt{|}
        \item \textbf{逻辑}：\texttt{\&\&} $\rightarrow$ \texttt{||}
        \item \textbf{三目}：\texttt{?:} \hfill \textit{(右到左)}
        \item \textbf{赋值}：\texttt{=} \texttt{+=} \texttt{-=} ... \hfill \textit{(右到左)}
        \item \textbf{逗号}：\texttt{,}
    \end{enumerate}
    \begin{tipbox}[title=优先级口诀与陷阱]
        口诀：单目 > 算术 > 关系 > 逻辑 > 三目 > 赋值
        \par\smallskip\hrule\smallskip\par
        \textbf{注意1}：关系运算符中，\textbf{大小比较} (\texttt{> <}) 高于 \textbf{相等比较} (\texttt{== !=})。
        \\ 例如：\texttt{a > b == c} 等价于 \texttt{(a > b) == c}。
        \par\smallskip\hrule\smallskip\par
        \textbf{注意2}：\textbf{右到左}结合性（单目、三目、赋值）。
        \\ \textbf{赋值}：\texttt{a = b = c = 0} $\rightarrow$ \texttt{a = (b = (c = 0))}。
        \\ \textbf{单目}：\texttt{*p++} $\rightarrow$ \texttt{*(p++)}（先取值后自增，再解引用原地址）。
    \end{tipbox}
\end{itemize}

\subsection{输入输出 (Input/Output)}
\begin{itemize}
    \item \textbf{格式化输出 (printf)}：
    \begin{itemize}
        \item \texttt{\%d} (整型), \texttt{\%f} (浮点), \texttt{\%c} (字符), \texttt{\%s} (字符串), \texttt{\%p} (指针地址)。
        \item \textbf{精度控制}：\texttt{\%.2f} (保留2位小数)，\texttt{\%5d} (宽5位右对齐)，\texttt{\%-5d} (宽5位左对齐)。
    \end{itemize}
    \item \textbf{格式化输入 (scanf)}：
    \begin{itemize}
        \item \textbf{地址符}：除了字符串数组名外，变量前必须加 \texttt{\&}。
        \item \textbf{分隔符}：若格式串为 \texttt{"\%d,\%d"}，输入时必须加逗号。
        \item \textbf{缓冲区残留}：混合输入字符和数字时，回车符可能被 \texttt{\%c} 读取。建议在 \texttt{\%c} 前加空格：\texttt{scanf(" \%c", \&ch);}。
    \end{itemize}
    \item \textbf{字符输入输出}：
    \begin{itemize}
        \item \texttt{getchar()} / \texttt{putchar()}：一次只处理一个字符。
    \end{itemize}
\end{itemize}

\subsection{常量与宏定义}
\begin{itemize}
    \item \textbf{const 常量}：
    \begin{itemize}
        \item \texttt{const int MAX = 100;}
        \item 有类型检查，占用内存（通常在只读区）。
    \end{itemize}
    \item \textbf{宏定义 (\#define)}：
    \begin{itemize}
        \item \texttt{\#define MAX 100}
        \item 预处理阶段文本替换，无类型检查，不占内存。
        \item \textbf{注意}：宏定义末尾不要加分号。
        \item \textbf{真题陷阱（带参宏）}：纯文本替换，不自动加括号。
        \begin{itemize}
            \item \textbf{错误}：\texttt{\#define S(x) x*x} $\rightarrow$ \texttt{S(a+b)} 变为 \texttt{a+b*a+b}。
            \item \textbf{正确}：\texttt{\#define S(x) ((x)*(x))} $\rightarrow$ \texttt{((a+b)*(a+b))}。
        \end{itemize}
    \end{itemize}
\end{itemize}

\subsection{易错点}
\begin{itemize}
    \item \textbf{除法陷阱}：\texttt{5 / 2 = 2}, \texttt{5.0 / 2 = 2.5}。
    \item \textbf{逻辑短路}：\texttt{A || B} (A真B不执行), \texttt{A \&\& B} (A假B不执行)。
    \item \textbf{浮点比较}：严禁 \texttt{f == 0.0}，应用 \texttt{fabs(f) < 1e-6}。
    \item \textbf{赋值与相等}：\texttt{if(a=2)} 永远为真，且 \texttt{a} 变为 2。这是极高频考点。
\end{itemize}

\section{控制结构}
\label{sec:control_structures}

\begin{itemize}
    \item \textbf{Switch-Case}：\texttt{case} 后必须是整型或字符型常量。注意\textbf{穿透性}（漏写 \texttt{break}）。
    \item \textbf{循环}：\texttt{continue} 跳过本次剩余，\texttt{break} 跳出循环。
    \item \textbf{易错点}：悬挂 \texttt{else}（与最近的 \texttt{if} 匹配）；分号错误（\texttt{if(x); ...}）。
    \item \textbf{真题模式}：利用 \texttt{do...while} 计算泰勒级数（如 \texttt{sin(x)}），注意循环终止条件通常是 \texttt{fabs(term) < 1e-5}。
\end{itemize}

\section{数组与字符串}

\subsection{基本概念}
\begin{itemize}
    \item \textbf{定义}：字符串是以空字符 \texttt{'\textbackslash 0'} 结尾的字符数组。
    \item \textbf{初始化}：
    \begin{itemize}
        \item \texttt{char s[] = "hello";} (自动加 \texttt{'\textbackslash 0'}，长度为6，内容可变)
        \item \texttt{char *p = "hello";} (字符串常量，通常存储在只读数据区，\textbf{不可修改})
    \end{itemize}
    \item \textbf{sizeof vs strlen}：
    \begin{itemize}
        \item \texttt{sizeof}：计算数组占用的总字节数（含 \texttt{'\textbackslash 0'}）。
        \item \texttt{strlen}：计算字符串的有效长度（不含 \texttt{'\textbackslash 0'}）。
        \item \textbf{真题陷阱}：\texttt{strlen("abc\textbackslash 0def")} 结果是 3（遇到第一个 \texttt{\textbackslash 0} 停止）。
    \end{itemize}
\end{itemize}

\subsection{常用字符串函数 (<string.h>)}
\label{sec:string_funcs}
\begin{itemize}
    \item \textbf{复制 (Copy)}：
    \begin{itemize}
        \item \texttt{strcpy(dest, src)}：将 \texttt{src} 复制到 \texttt{dest}。
        \item \textbf{危险}：若 \texttt{dest} 空间不足，会发生缓冲区溢出。
        \item \texttt{strncpy(dest, src, n)}：最多复制 \texttt{n} 个字符。注意：若 \texttt{src} 长度 $> n$，则 \texttt{dest} 不会自动添加 \texttt{'\textbackslash 0'}，需手动添加。
    \end{itemize}
    
    \item \textbf{拼接 (Concatenate)}：
    \begin{itemize}
        \item \texttt{strcat(dest, src)}：将 \texttt{src} 拼接到 \texttt{dest} 后面。
        \item \textbf{要求}：\texttt{dest} 必须有足够的剩余空间，且必须以 \texttt{'\textbackslash 0'} 结尾。
        \item \texttt{strncat(dest, src, n)}：拼接前 \texttt{n} 个字符，会自动添加 \texttt{'\textbackslash 0'}。
    \end{itemize}
    
    \item \textbf{比较 (Compare)}：
    \begin{itemize}
        \item \texttt{strcmp(s1, s2)}：按字典序比较。
        \item \textbf{返回值}：0 (相等), >0 (s1大), <0 (s1小)。
        \item \textbf{陷阱}：严禁使用 \texttt{if(s1 == s2)} 比较内容（这是比较地址！）。
    \end{itemize}
    
    \item \textbf{查找 (Search)}：
    \begin{itemize}
        \item \texttt{strchr(s, c)}：查找字符 \texttt{c} 首次出现的位置（返回指针）。
        \item \texttt{strstr(s1, s2)}：查找子串 \texttt{s2} 首次出现的位置。
    \end{itemize}
\end{itemize}

\subsection{其他重要函数}
\begin{itemize}
    \item \texttt{sprintf(buf, "\%d", 123)}：将数据格式化输出到字符串中（常用于数字转字符串）。
    \item \texttt{atoi(s) / atof(s)}：字符串转整数/浮点数 (\texttt{<stdlib.h>})。
\end{itemize}

\section{指针 (核心难点)}

\subsection{指针基础与运算}
\label{sec:pointer_basics}
\begin{itemize}
    \item \textbf{定义与解引用}：
    \begin{itemize}
        \item \texttt{int *p = \&a;}：\texttt{p} 存储 \texttt{a} 的地址。
        \item \texttt{*p}：访问 \texttt{p} 指向的内存单元（即 \texttt{a}）。
        \item \textbf{口诀}：\texttt{\&} 取地址，\texttt{*} 取值。
    \end{itemize}
    \item \textbf{指针运算}：
    \begin{itemize}
        \item \texttt{p + 1}：\textbf{不是地址加1}，而是地址增加 \texttt{sizeof(类型)}。
        \item \texttt{p1 - p2}：返回两指针之间的\textbf{元素个数}（不是字节数）。
        \item \textbf{示例}：若 \texttt{int a[5]}，\texttt{p=a}，则 \texttt{*(p+2)} 等价于 \texttt{a[2]}。
    \end{itemize}
\end{itemize}

\subsection{指针与数组的深度辨析}
\label{sec:pointer_array}
\begin{itemize}
    \item \textbf{访问方式等价性}：
    \begin{itemize}
        \item \texttt{a[i]} $\iff$ \texttt{*(a+i)} $\iff$ \texttt{*(p+i)} $\iff$ \texttt{p[i]}。
    \end{itemize}
    \item \textbf{本质区别}：
    \begin{itemize}
        \item \textbf{数组名} \texttt{a} 是\textbf{常量指针}，不可修改（\texttt{a++} 非法）。
        \item \textbf{指针变量} \texttt{p} 是变量，可以修改指向（\texttt{p++} 合法）。
    \end{itemize}
    \item \textbf{二维数组指针}：
    \begin{itemize}
        \item \texttt{int a[3][4];}
        \item \texttt{a}：指向首行（类型 \texttt{int (*)[4]}）。
        \item \texttt{a[0]}：指向首行首元素（类型 \texttt{int *}）。
        \item \texttt{a[i][j]} $\iff$ \texttt{*(*(a+i)+j)}。
    \end{itemize}
\end{itemize}

\subsection{复杂指针声明解析}
\label{sec:complex_pointers}
\begin{itemize}
    \item \textbf{const修饰}（“左定值，右定向”）：
    \begin{itemize}
        \item \texttt{const int *p}：\texttt{*p} (内容) 不可变，\texttt{p} (指向) 可变。
        \item \texttt{int * const p}：\texttt{p} (指向) 不可变，\texttt{*p} (内容) 可变。
    \end{itemize}
    \item \textbf{指针数组 vs 数组指针}：
    \begin{itemize}
        \item \texttt{int *p[10]}：\textbf{数组}，存了10个指针。
        \item \texttt{int (*p)[10]}：\textbf{指针}，指向长度为10的数组。
    \end{itemize}
    \item \textbf{函数指针}：
    \begin{itemize}
        \item 定义：\texttt{int (*func)(int, int);}
        \item 用途：回调函数（如 \texttt{qsort} 的比较函数）。
    \end{itemize}
\end{itemize}

\subsection{简答题预测}
\begin{description}
    \item[Q: 什么是指针？为什么要使用指针？] \hfill \\
    A: 指针是存储内存地址的变量。作用：1. 直接操作内存；2. 函数间传递大数据（避免拷贝）；3. 动态内存分配；4. 操作复杂数据结构。
\end{description}

\section{函数与递归}
\label{sec:functions}

\begin{itemize}
    \item \textbf{传值 vs 传址}：修改实参需传指针。数组传参退化为指针。
    \item \textbf{递归 (Recursion)}：
    \begin{itemize}
        \item \textbf{要素}：基准情况 (Base Case) + 递归步骤 (Recursive Step)。
        \item \textbf{优缺点}：代码简洁但开销大（栈空间）。
    \end{itemize}
    \item \textbf{存储类别}：
    \begin{itemize}
        \item \texttt{auto}：栈内存。
        \item \texttt{static}：延长局部变量生命周期；限制全局变量作用域。
        \item \textbf{真题考点}：局部变量与全局变量同名时，局部变量优先（屏蔽全局变量）。
    \end{itemize}
\end{itemize}

\section{结构体、共用体与枚举}

\subsection{核心考点}
\begin{itemize}
    \item \textbf{结构体 (Struct)}：
    \begin{itemize}
        \item 成员拥有独立内存。
        \item \textbf{内存对齐}：总大小 $\ge$ 成员之和。
        \item \textbf{真题考点}：\texttt{typedef struct \{...\} Name;} 定义别名，方便使用。
    \end{itemize}
    \item \textbf{共用体 (Union)}：
    \begin{itemize}
        \item 成员共享同一块内存。
        \item 总大小 = 最大成员大小。
    \end{itemize}
    \item \textbf{枚举 (Enum)}：默认从0开始递增。
\end{itemize}

\subsection{简答题预测}
\begin{description}
    \item[Q: 简述结构体和共用体的区别。] \hfill \\
    A: 1. 内存：结构体成员独立，共用体成员共享。2. 大小：结构体累加（含对齐），共用体取最大。3. 赋值：共用体同一时间只能存一个值。
\end{description}

\section{动态内存管理 (Dynamic Memory)}

\subsection{栈 (Stack) 与 堆 (Heap) 的区别}
\begin{itemize}
    \item \textbf{栈 (Stack)}：
    \begin{itemize}
        \item 由编译器自动分配释放。
        \item 存放局部变量、函数参数。
        \item 空间较小，速度快。
    \end{itemize}
    \item \textbf{堆 (Heap)}：
    \begin{itemize}
        \item 由程序员手动分配 (\texttt{malloc}) 和释放 (\texttt{free})。
        \item 空间大，但容易产生内存泄漏或碎片。
    \end{itemize}
\end{itemize}

\subsection{核心函数详解}
\begin{itemize}
    \item \textbf{malloc}：\texttt{void *malloc(size\_t size);}
    \begin{itemize}
        \item 分配指定字节数的内存。
        \item \textbf{内容未初始化}（随机值）。
        \item 返回 \texttt{void*}，通常需要强制类型转换（C语言中可省略，但建议写上）。
        \item \textbf{示例}：\texttt{int *p = (int*)malloc(n * sizeof(int));}
    \end{itemize}
    \item \textbf{calloc}：\texttt{void *calloc(size\_t n, size\_t size);}
    \begin{itemize}
        \item 分配 $n \times size$ 字节。
        \item \textbf{自动初始化为 0}。
    \end{itemize}
    \item \textbf{realloc}：\texttt{void *realloc(void *ptr, size\_t new\_size);}
    \begin{itemize}
        \item 调整已分配内存的大小。
        \item 可能移动内存块位置，必须接收返回值更新指针。
    \end{itemize}
    \item \textbf{free}：\texttt{void free(void *ptr);}
    \begin{itemize}
        \item 释放内存，归还给系统。
        \item \textbf{注意}：释放后指针变量本身的值未变（悬空指针），建议立即置为 \texttt{NULL}。
    \end{itemize}
\end{itemize}

\subsection{常见错误与最佳实践}
\begin{itemize}
    \item \textbf{忘记检查 NULL}：
    \begin{lstlisting}
    int *p = (int*)malloc(1000 * sizeof(int));
    if (p == NULL) {
        // 处理分配失败
        exit(1);
    }
    \end{lstlisting}
    \item \textbf{内存泄漏 (Memory Leak)}：申请了内存但没有 \texttt{free}，导致程序占用内存越来越多。
    \item \textbf{悬空指针 (Dangling Pointer)}：使用已经 \texttt{free} 过的指针。
    \item \textbf{重复释放}：对同一个指针 \texttt{free} 两次（会导致崩溃）。
\end{itemize}

\section{链表与队列 (Linked List \& Queue)}

\subsection{单向链表 (Singly Linked List)}
链表是一种物理存储单元上非连续、非顺序的存储结构，数据元素的逻辑顺序是通过链表中的指针链接次序实现的。

\subsubsection{节点定义}
\begin{lstlisting}[style=mystyle]
typedef struct Node {
    int data;               // 数据域
    struct Node *next;      // 指针域，指向下一个节点
} Node;
\end{lstlisting}

\subsubsection{链表的核心操作}
\begin{enumerate}
    \item \textbf{创建新节点}
    \begin{lstlisting}[style=mystyle]
Node* createNode(int value) {
    Node* newNode = (Node*)malloc(sizeof(Node));
    if (newNode == NULL) {
        printf("Memory allocation failed!\n");
        exit(1);
    }
    newNode->data = value;
    newNode->next = NULL;
    return newNode;
}
    \end{lstlisting}

    \item \textbf{头插法 (Insert at Head)}
    \begin{lstlisting}[style=mystyle]
void insertHead(Node** head, int value) {
    Node* newNode = createNode(value);
    newNode->next = *head; // 新节点的 next 指向当前头节点
    *head = newNode;       // 更新头指针指向新节点
}
    \end{lstlisting}

    \item \textbf{尾插法 (Insert at Tail)}
    \begin{lstlisting}[style=mystyle]
void insertTail(Node** head, int value) {
    Node* newNode = createNode(value);
    if (*head == NULL) {
        *head = newNode;
        return;
    }
    Node* temp = *head;
    while (temp->next != NULL) { // 遍历到最后一个节点
        temp = temp->next;
    }
    temp->next = newNode;
}
    \end{lstlisting}

    \item \textbf{链表遍历}
    \begin{lstlisting}[style=mystyle]
void printList(Node* head) {
    Node* current = head;
    while (current != NULL) {
        printf("%d -> ", current->data);
        current = current->next;
    }
    printf("NULL\n");
}
    \end{lstlisting}
\end{enumerate}

\subsection{队列 (Queue)}
队列是一种特殊的线性表，只允许在表的前端（front）进行删除操作，在表的后端（rear）进行插入操作。遵循 \textbf{FIFO (First In First Out, 先进先出)} 原则。

\subsubsection{队列的结构体定义 (链式队列)}
使用链表实现队列可以避免数组大小固定的限制。

\begin{lstlisting}[style=mystyle]
typedef struct {
    Node *front; // 队头指针，用于出队
    Node *rear;  // 队尾指针，用于入队
} Queue;

// 初始化队列
void initQueue(Queue *q) {
    q->front = q->rear = NULL;
}
\end{lstlisting}

\subsubsection{入队与出队}
\begin{itemize}
    \item \textbf{入队 (Enqueue)}：在 \texttt{rear} 后添加新节点。
    \begin{lstlisting}[style=mystyle]
void enqueue(Queue *q, int value) {
    Node* newNode = createNode(value);
    if (q->rear == NULL) { // 队列为空
        q->front = q->rear = newNode;
        return;
    }
    q->rear->next = newNode; // 旧队尾指向新节点
    q->rear = newNode;       // 更新队尾指针
}
    \end{lstlisting}

    \item \textbf{出队 (Dequeue)}：移除 \texttt{front} 指向的节点。
    \begin{lstlisting}[style=mystyle]
int dequeue(Queue *q) {
    if (q->front == NULL) return -1; // 队列为空

    Node *temp = q->front;
    int data = temp->data;
    
    q->front = q->front->next; // 队头指针后移
    
    if (q->front == NULL) { // 如果出队后队列变空，rear也要置空
        q->rear = NULL;
    }
    
    free(temp); // 释放节点内存
    return data;
}
    \end{lstlisting}
\end{itemize}

\subsection{易错点}
\begin{itemize}
    \item \textbf{内存泄漏}：删除节点或销毁链表时，记得 \texttt{free} 内存。
    \item \textbf{空链表处理}：在插入、删除操作时，要考虑头指针为 \texttt{NULL} 的情况。
    \item \textbf{断链}：插入节点时，先连后断。
\end{itemize}

\section{文件操作详解 (File I/O)}

\subsection{文件打开与关闭}
\begin{itemize}
    \item \textbf{打开文件}：\texttt{FILE *fp = fopen("filename", "mode");}
    \item \textbf{检查是否成功}：\textbf{必须}检查 \texttt{fp} 是否为 \texttt{NULL}。
    \begin{lstlisting}
    if ((fp = fopen("data.txt", "r")) == NULL) {
        printf("Cannot open file.\n");
        exit(1);
    }
    \end{lstlisting}
    \item \textbf{关闭文件}：\texttt{fclose(fp);} 也就是保存文件，释放缓冲区。
\end{itemize}

\subsection{文件打开模式 (Mode) 深度解析}
\begin{itemize}
    \item \textbf{文本模式 (Text Mode)}：
    \begin{itemize}
        \item \texttt{"r"} (read)：\textbf{只读}。
        \begin{itemize}
            \item 文件\textbf{必须存在}，否则返回 \texttt{NULL}。
            \item 指针指向文件开头。
        \end{itemize}
        \item \texttt{"w"} (write)：\textbf{只写}。
        \begin{itemize}
            \item 文件不存在则\textbf{创建}。
            \item 文件存在则\textbf{清空} (Truncate to zero length) —— \textbf{极度危险，慎用！}
            \item 指针指向文件开头。
        \end{itemize}
        \item \texttt{"a"} (append)：\textbf{追加}。
        \begin{itemize}
            \item 文件不存在则\textbf{创建}。
            \item 文件存在则\textbf{保留原内容}，指针指向文件\textbf{末尾}。
            \item 只能在末尾写，不能修改中间内容。
        \end{itemize}
        \item \texttt{"r+"}：\textbf{读写} (覆盖模式)。
        \begin{itemize}
            \item 文件\textbf{必须存在}。
            \item 指针指向开头。可以读，也可以从头开始覆盖写。
        \end{itemize}
        \item \texttt{"w+"}：\textbf{读写} (清空模式)。
        \begin{itemize}
            \item 文件不存在则创建，存在则\textbf{清空}。
            \item 既能读也能写，但刚打开时文件是空的。
        \end{itemize}
        \item \texttt{"a+"}：\textbf{读写追加}。
        \begin{itemize}
            \item 文件不存在则创建，存在则保留。
            \item \textbf{读}：指针默认在开头（可以读旧数据）。
            \item \textbf{写}：永远在末尾追加（无论指针在哪）。
        \end{itemize}
    \end{itemize}

    \item \textbf{二进制模式 (Binary Mode)}：
    \begin{itemize}
        \item 标志：在模式后加 \texttt{b}，如 \texttt{"rb"}, \texttt{"wb"}, \texttt{"ab+"}。
        \item \textbf{核心区别}：换行符处理。
        \begin{itemize}
            \item \textbf{文本模式}：在 Windows 下，写入 \texttt{\textbackslash n} 会自动转换为 \texttt{\textbackslash r\textbackslash n} (CRLF)；读取 \texttt{\textbackslash r\textbackslash n} 会自动转换为 \texttt{\textbackslash n}。Linux 下无此转换。
            \item \textbf{二进制模式}：不做任何转换，读写原始字节。
        \end{itemize}
        \item \textbf{适用场景}：图片、音频、视频、结构体数据块 (\texttt{fread}/\texttt{fwrite}) 必须用二进制模式，否则数据会损坏。
    \end{itemize}
\end{itemize}

\subsection{读写函数辨析}
\begin{itemize}
    \item \textbf{字符读写}：
    \begin{itemize}
        \item \texttt{ch = fgetc(fp)}：读一个字符。结束返回 \texttt{EOF}。
        \item \texttt{fputc(ch, fp)}：写一个字符。
    \end{itemize}
    \item \textbf{字符串读写}：
    \begin{itemize}
        \item \texttt{fgets(str, n, fp)}：读取一行（最多 $n-1$ 个字符），\textbf{保留换行符}，自动加 \texttt{\textbackslash 0}。
        \item \texttt{fputs(str, fp)}：写入字符串，不会自动加换行符。
    \end{itemize}
    \item \textbf{格式化读写}（文本文件）：
    \begin{itemize}
        \item \texttt{fscanf(fp, "\%d", \&num)}：从文件读。
        \item \texttt{fprintf(fp, "Result: \%d", num)}：写入文件。
    \end{itemize}
    \item \textbf{数据块读写}（二进制文件）：
    \begin{itemize}
        \item \texttt{fread(buffer, size, count, fp)}：从文件读 count 个 size 大小的块到 buffer。
        \item \texttt{fwrite(buffer, size, count, fp)}：将 buffer 中 count 个 size 大小的块写入文件。
        \item \textbf{返回值}：成功读/写的块数 (\texttt{count})。
    \end{itemize}
\end{itemize}

\subsection{文件定位与随机读写}
\begin{itemize}
    \item \texttt{rewind(fp)}：文件指针回到开头。
    \item \texttt{ftell(fp)}：返回当前文件指针位置（字节偏移量）。
    \item \texttt{fseek(fp, offset, origin)}：移动指针。
    \begin{itemize}
        \item \texttt{SEEK\_SET} (0)：文件头。
        \item \texttt{SEEK\_CUR} (1)：当前位置。
        \item \texttt{SEEK\_END} (2)：文件尾。
    \end{itemize}
\end{itemize}

\subsection{文件操作常见考点与易错点}
\begin{itemize}
    \item \textbf{feof 的误用 (经典考点)}：
    \begin{itemize}
        \item \textbf{错误写法}：\texttt{while(!feof(fp))}。
        \item \textbf{原因}：\texttt{feof} 只有在\textbf{读取操作尝试读取文件尾之后}才会返回真，而不是到达文件尾就立即返回真。这会导致循环多执行一次（最后一次读取失败但循环体仍执行）。
        \item \textbf{正确写法}：利用读取函数的返回值作为循环条件。
        \begin{lstlisting}
        // 文本读整数
        while (fscanf(fp, "%d", &n) == 1) { ... }
        // 文本读行
        while (fgets(buf, 100, fp) != NULL) { ... }
        // 二进制读块
        while (fread(&data, sizeof(data), 1, fp) == 1) { ... }
        \end{lstlisting}
    \end{itemize}
    
    \item \textbf{文本文件 vs 二进制文件存储大小}：
    \begin{itemize}
        \item 考题常问：整数 \texttt{12345} 存入文件占用多少字节？
        \item \textbf{文本模式}：按字符存，占用 \textbf{5个字节} ('1','2','3','4','5')。
        \item \textbf{二进制模式}：按内存原样存 (int)，通常占用 \textbf{4个字节}。
    \end{itemize}

    \item \textbf{标准流 (Standard Streams)}：
    \begin{itemize}
        \item 系统自动打开的三个文件指针：
        \item \texttt{stdin} (标准输入，对应键盘)
        \item \texttt{stdout} (标准输出，对应屏幕)
        \item \texttt{stderr} (标准错误，对应屏幕，无缓冲)
        \item \textbf{技巧}：\texttt{fprintf(stdout, ...)} 等价于 \texttt{printf(...)}。
    \end{itemize}

    \item \textbf{文件打开失败检测}：
    \begin{itemize}
        \item 必须检查 \texttt{fopen} 的返回值是否为 \texttt{NULL}。
        \begin{lstlisting}
        FILE *fp = fopen("data.txt", "r");
        if (fp == NULL) {
            perror("File open error"); // 输出错误原因
            return -1;
        }
        \end{lstlisting}
    \end{itemize}

    \item \textbf{读写同步 (Update Mode)}：
    \begin{itemize}
        \item 在 \texttt{r+}, \texttt{w+} 等模式下，\textbf{读}与\textbf{写}操作之间必须插入 \texttt{fseek}, \texttt{fsetpos} 或 \texttt{fflush}，否则文件指针位置会混乱导致数据错误。
    \end{itemize}
    
    \item \textbf{路径分隔符}：
    \begin{itemize}
        \item Windows 路径 \texttt{C:\textbackslash test.txt} 在 C 字符串中必须转义：\texttt{"C:\textbackslash\textbackslash test.txt"} 或使用正斜杠 \texttt{"C:/test.txt"}。
    \end{itemize}
\end{itemize}

\subsection{预处理指令补充}
\begin{itemize}
    \item \textbf{Include}：\texttt{<>} 查系统库，\texttt{""} 查当前目录。
\end{itemize}

\section{核心算法详解与易错点深度剖析}

\subsection{排序算法 (Sorting)}

\subsubsection{冒泡排序 (Bubble Sort)}
\textbf{核心思想}：像水底的气泡一样，每一轮通过两两比较和交换，将当前未排序部分的最大值“浮”到数组的最右端。

\begin{lstlisting}
void bubble_sort(int arr[], int n) {
    int i, j, temp;
    // 外层循环：控制排序轮数，共需 n-1 轮
    for (i = 0; i < n - 1; i++) {
        // 内层循环：负责两两比较
        // 易错点1：边界是 n-1-i，因为最后 i 个元素已经排好序了
        for (j = 0; j < n - 1 - i; j++) {
            if (arr[j] > arr[j + 1]) { // 升序用 >，降序用 <
                temp = arr[j];
                arr[j] = arr[j + 1];
                arr[j + 1] = temp;
            }
        }
    }
}
\end{lstlisting}

\textbf{🔴 易错点与考点}：
\begin{itemize}
    \item \textbf{内层循环边界}：\texttt{j < n - 1 - i}。
    \begin{itemize}
        \item 如果写成 \texttt{j < n - 1}，程序也能跑，但效率低。
        \item 如果写成 \texttt{j < n}，在访问 \texttt{arr[j+1]} 时会\textbf{数组越界}。
    \end{itemize}
    \item \textbf{交换位置}：交换是在\textbf{内层循环}的 \texttt{if} 语句中发生的。
    \item \textbf{稳定性}：冒泡排序是\textbf{稳定}的。
    \item \textbf{时间复杂度}：平均和最坏情况均为 $O(n^2)$。
    \begin{itemize}
        \item \textbf{推导}：外层循环执行 $n-1$ 次。内层循环第 1 轮比较 $n-1$ 次，第 2 轮比较 $n-2$ 次……直到最后 1 次。
        \item 总比较次数 = $(n-1) + (n-2) + \dots + 1 = \frac{n(n-1)}{2}$，属于 $O(n^2)$ 级别。
    \end{itemize}
\end{itemize}

\subsubsection{选择排序 (Selection Sort)}
\textbf{核心思想}：每一轮从未排序区域中找到\textbf{最小值}的下标，然后将其与未排序区域的第一个元素交换。

\begin{lstlisting}
void selection_sort(int arr[], int n) {
    int i, j, min_idx, temp;
    // 外层循环：控制当前要确定的位置，从 0 到 n-2
    for (i = 0; i < n - 1; i++) {
        min_idx = i; // 假设当前位置 i 就是最小值
        
        // 内层循环：在剩余元素中找真正的最小值下标
        for (j = i + 1; j < n; j++) {
            if (arr[j] < arr[min_idx]) {
                min_idx = j; // 更新最小值下标，注意这里不交换！
            }
        }
        
        // 易错点2：交换发生在内层循环结束后
        if (min_idx != i) {
            temp = arr[i];
            arr[i] = arr[min_idx];
            arr[min_idx] = temp;
        }
    }
}
\end{lstlisting}

\textbf{🔴 易错点与考点}：
\begin{itemize}
    \item \textbf{交换时机}：千万不要在内层循环里交换！内层循环只负责\textbf{找下标} (\texttt{min\_idx})。
    \item \textbf{比较对象}：内层循环是 \texttt{if (arr[j] < arr[min\_idx])}，而不是 \texttt{arr[j] < arr[i]}。
    \item \textbf{稳定性}：选择排序是\textbf{不稳定}的。
    \item \textbf{时间复杂度}：无论数组是否有序，均为 $O(n^2)$。
    \begin{itemize}
        \item \textbf{推导}：无论数据如何，内层循环的比较次数固定。
        \item 第 1 轮找最小值比较 $n-1$ 次，第 2 轮比较 $n-2$ 次……总次数恒为 $\frac{n(n-1)}{2}$。
    \end{itemize}
\end{itemize}

\subsubsection{插入排序 (Insertion Sort)}
\textbf{核心思想}：像打扑克牌理牌一样，将新抓到的牌（\texttt{key}）插入到左手已经排好序的牌堆中的正确位置。

\begin{lstlisting}
void insertion_sort(int arr[], int n) {
    int i, j, key;
    // 从第2个元素开始（下标1），因为第1个元素默认有序
    for (i = 1; i < n; i++) {
        key = arr[i]; // 待插入的牌
        j = i - 1;    // 左边已排序区域的最后一个元素
        
        // 易错点3：while 循环条件
        // j >= 0 保证不越界
        // arr[j] > key 保证找到比 key 大的元素往后挪
        while (j >= 0 && arr[j] > key) {
            arr[j + 1] = arr[j]; // 元素后移
            j--; // 继续向前找
        }
        arr[j + 1] = key; // 插入到正确位置
    }
}
\end{lstlisting}

\textbf{🔴 易错点与考点}：
\begin{itemize}
    \item \textbf{后移操作}：\texttt{arr[j + 1] = arr[j]} 是覆盖操作，所以必须先保存 \texttt{key}。
    \item \textbf{插入位置}：循环结束时，\texttt{key} 应该放在 \texttt{j + 1} 的位置。
    \item \textbf{稳定性}：插入排序是\textbf{稳定}的。
    \item \textbf{时间复杂度}：
    \begin{itemize}
        \item \textbf{最坏情况} ($O(n^2)$)：数组完全逆序。插入第 $i$ 个元素时，需要向前比较 $i$ 次。总次数 $\sum_{i=1}^{n-1} i = \frac{n(n-1)}{2}$。
        \item \textbf{最好情况} ($O(n)$)：数组已有序。内层 \texttt{while} 循环条件一开始就不满足（\texttt{arr[j] < key}），每次只比较 1 次，不移动。总共比较 $n-1$ 次。
    \end{itemize}
\end{itemize}

\subsection{查找算法 (Searching)}

\subsubsection{二分查找 (Binary Search)}
\textbf{前提条件}：数组必须是\textbf{有序}的。

\begin{lstlisting}
int binary_search(int arr[], int n, int key) {
    int low = 0, high = n - 1;
    int mid;
    
    while (low <= high) { // 易错点4：是 <= 不是 <
        // 易错点5：防溢出写法
        mid = low + (high - low) / 2; 
        
        if (arr[mid] == key) {
            return mid; // 找到了
        } else if (arr[mid] < key) {
            low = mid + 1; // 在右半边，low 必须 +1
        } else {
            high = mid - 1; // 在左半边，high 必须 -1
        }
    }
    return -1; // 没找到
}
\end{lstlisting}

\textbf{🔴 易错点与考点}：
\begin{itemize}
    \item \textbf{循环条件}：\texttt{while (low <= high)}。如果写成 \texttt{<}，会漏掉边界情况。
    \item \textbf{死循环陷阱}：更新边界时必须是 \texttt{mid + 1} 或 \texttt{mid - 1}。
    \item \textbf{Mid 计算}：\texttt{(low + high) / 2} 可能溢出，推荐 \texttt{low + (high - low) / 2}。
\end{itemize}

\subsection{数学与逻辑算法}

\subsubsection{辗转相除法 (GCD)}
\label{sec:gcd}
\textbf{核心思想}：\texttt{gcd(a, b) = gcd(b, a \% b)}。

\begin{lstlisting}
// 递归版
int gcd(int a, int b) {
    if (b == 0) return a;
    return gcd(b, a % b);
}
\end{lstlisting}

\subsubsection{素数判断 (Prime Check)}
\textbf{核心思想}：除了 1 和它本身外，不能被其他数整除。

\begin{lstlisting}
#include <math.h> // 需要用到 sqrt

int is_prime(int n) {
    if (n <= 1) return 0; // 易错点6：1 不是素数
    
    // 优化：只需要判断到 sqrt(n)
    int limit = (int)sqrt(n);
    for (int i = 2; i <= limit; i++) {
        if (n % i == 0) {
            return 0; // 发现因子，不是素数
        }
    }
    return 1; // 是素数
}
\end{lstlisting}

\textbf{🔴 易错点与考点}：
\begin{itemize}
    \item \textbf{特判 1}：\texttt{1} 既不是质数也不是合数，必须单独判断返回 0。
    \item \textbf{循环终止条件}：\texttt{i * i <= n} 或 \texttt{i <= sqrt(n)}。
\end{itemize}

\subsubsection{常见数学公式实现}
\label{sec:math_formulas}
\begin{itemize}
    \item \textbf{海伦公式 (Heron's Formula)}：已知三角形三边 $a, b, c$，求面积。
    \[ p = \frac{a+b+c}{2}, \quad S = \sqrt{p(p-a)(p-b)(p-c)} \]
    \begin{lstlisting}
#include <math.h>
double area(double a, double b, double c) {
    double p = (a + b + c) / 2.0;
    return sqrt(p * (p - a) * (p - b) * (p - c));
}
    \end{lstlisting}
    \item \textbf{泰勒级数求 sin(x)}：$\sin(x) = x - \frac{x^3}{3!} + \frac{x^5}{5!} - \dots$
    \begin{lstlisting}
#include <math.h>
double my_sin(double x) {
    double sum = 0, term = x;
    int n = 1;
    while (fabs(term) >= 1e-5) { // 精度控制
        sum += term;
        term = -term * x * x / ((n + 1) * (n + 2)); // 递推公式
        n += 2;
    }
    return sum;
}
    \end{lstlisting}
\end{itemize}

\subsection{字符串处理算法}

\subsubsection{字符串逆序}
\textbf{核心思想}：双指针法，首尾交换，向中间逼近。

\begin{lstlisting}
void reverse_string(char str[]) {
    int len = strlen(str);
    int i = 0;          // 头指针
    int j = len - 1;    // 尾指针
    char temp;
    
    while (i < j) {
        temp = str[i];
        str[i] = str[j];
        str[j] = temp;
        i++;
        j--;
    }
}
\end{lstlisting}

\textbf{🔴 易错点与考点}：
\begin{itemize}
    \item \textbf{尾指针位置}：\texttt{len - 1}。
    \item \textbf{循环条件}：\texttt{i < j}。
    \item \textbf{不要动 \textbackslash 0}：\texttt{\textbackslash 0} 是字符串结束符，不参与交换。
\end{itemize}

\subsection{考试手写代码“救命”技巧}
\begin{enumerate}
    \item \textbf{头文件}：用到 \texttt{printf} 写 \texttt{\#include <stdio.h>}，用到 \texttt{malloc} 写 \texttt{\#include <stdlib.h>}，用到 \texttt{strlen} 写 \texttt{\#include <string.h>}，用到 \texttt{sqrt} 写 \texttt{\#include <math.h>}。
    \item \textbf{变量初始化}：计数器 \texttt{count}、累加器 \texttt{sum} 一定要初始化为 0。
    \item \textbf{分号检查}：\texttt{if}、\texttt{for}、\texttt{while} 后面千万别手顺加分号！
    \item \textbf{返回值}：\texttt{main} 函数最后记得写 \texttt{return 0;}。
\end{enumerate}

\section{C语言易混淆概念辨析}

\begin{itemize}
    \item \textbf{赋值 vs 相等}：
    \begin{itemize}
        \item \texttt{=} 是赋值运算符（把右边的值给左边）。
        \item \texttt{==} 是关系运算符（判断左右两边是否相等）。
        \item \textbf{陷阱}：\texttt{if(x=5)} 永远为真，\texttt{x} 变为5；\texttt{if(x==5)} 才是判断。
    \end{itemize}

    \item \textbf{字符 vs 字符串}：
    \begin{itemize}
        \item \texttt{'a'}：字符常量，占1字节。
        \item \texttt{"a"}：字符串常量，包含 \texttt{'a'} 和 \texttt{'\textbackslash 0'}，占2字节。
    \end{itemize}

    \item \textbf{逻辑与 vs 按位与}：
    \begin{itemize}
        \item \texttt{\&\&} (逻辑与)：短路特性，结果为 0 或 1。
        \item \texttt{\&} (按位与)：按二进制位操作，如 \texttt{3 \& 5} (011 \& 101 = 001)。
    \end{itemize}

    \item \textbf{sizeof vs strlen}：
    \begin{itemize}
        \item \texttt{sizeof}：运算符，编译时计算，统计类型或变量占用的\textbf{内存字节数}（包含 \texttt{\textbackslash 0}）。
        \item \texttt{strlen}：函数，运行时计算，统计字符串的\textbf{有效字符长度}（遇到 \texttt{\textbackslash 0} 停止，不含 \texttt{\textbackslash 0}）。
    \end{itemize}

    \item \textbf{指针数组 vs 数组指针}：
    \begin{itemize}
        \item \texttt{int *p[10]}：\textbf{数组}，里面存了10个 \texttt{int*} 指针。
        \item \texttt{int (*p)[10]}：\textbf{指针}，指向一个包含10个 \texttt{int} 的数组。
    \end{itemize}

    \item \textbf{p++ vs (*p)++}：
    \begin{itemize}
        \item \texttt{p++}：指针向后移动一个单位（地址改变）。
        \item \texttt{(*p)++}：指针指向的值加1（地址不变，内容改变）。
    \end{itemize}

    \item \textbf{break vs continue}：
    \begin{itemize}
        \item \texttt{break}：彻底跳出当前循环（或 switch）。
        \item \texttt{continue}：跳过本次循环剩余语句，直接进入下一次循环判断。
    \end{itemize}
    
    \item \textbf{形参 vs 实参}：
    \begin{itemize}
        \item \textbf{单向传递}：实参的值传给形参，形参改变不影响实参（除非传指针）。
    \end{itemize}
\end{itemize}

\section{冷门与进阶考点拾遗}

\begin{itemize}
    \item \textbf{逗号运算符}：
    \begin{itemize}
        \item 优先级最低。
        \item 表达式的值是\textbf{最后一个表达式}的值。
        \item 例：\texttt{a = 3, 4, 5;} $\rightarrow$ \texttt{a} 为 3（赋值优先级高于逗号）。
        \item 例：\texttt{a = (3, 4, 5);} $\rightarrow$ \texttt{a} 为 5。
    \end{itemize}

    \item \textbf{宏定义陷阱}：
    \begin{itemize}
        \item 宏只是简单的文本替换，不进行计算。
        \item 例：\texttt{\#define S(r) r*r}。若调用 \texttt{S(a+b)}，展开为 \texttt{a+b*a+b}，而非 \texttt{(a+b)*(a+b)}。
        \item \textbf{对策}：宏定义中每个参数都要加括号，整体也要加括号。
    \end{itemize}

    \item \textbf{转义字符 (Escape Characters)}：
    \begin{itemize}
        \item 详见 \textbf{附录4：字符与字符串陷阱深度解析}。
        \item 核心考点：\texttt{\textbackslash ddd} (八进制) vs \texttt{\textbackslash xhh} (十六进制) 的长度限制与非法值判断。
    \end{itemize}

    \item \textbf{格式化输出细节}：
    \begin{itemize}
        \item \texttt{\%5d}：域宽为5，右对齐。
        \item \texttt{\%-5d}：域宽为5，左对齐。
        \item \texttt{\%05d}：域宽为5，不足补0。
        \item \texttt{\%5.2f}：总宽5，小数点后2位。
    \end{itemize}

    \item \textbf{运算符优先级口诀}：
    \begin{itemize}
        \item \textbf{单目} (!, ++, --, sizeof) > \textbf{算术} (*, /, \%) > \textbf{算术} (+, -) > \textbf{移位} (<<, >>) > \textbf{关系} (>, <) > \textbf{相等} (==, !=) > \textbf{位逻辑} (\&, \^{}, |) > \textbf{逻辑} (\&\&, ||) > \textbf{三目} (? :) > \textbf{赋值} (=) > \textbf{逗号} (,)。
        \item 记不住就加括号！
    \end{itemize}
    
    \item \textbf{typedef 与 \#define 的区别}：
    \begin{itemize}
        \item \texttt{\#define} 是预处理指令，仅做文本替换。
        \item \texttt{typedef} 是编译指令，为类型创建别名，有作用域检查。
        \item 例：\texttt{typedef int* ptr;} \texttt{ptr a, b;} $\rightarrow$ a, b 都是指针。
        \item 例：\texttt{\#define PTR int*} \texttt{PTR a, b;} $\rightarrow$ a 是指针，b 是整型（替换后为 \texttt{int * a, b;}）。
    \end{itemize}
\end{itemize}

\section{全卷易错点与陷阱大盘点}
\label{sec:common_mistakes}

本章节汇总了全卷最容易丢分的陷阱，建议考前反复查阅！

\subsection{语法与逻辑陷阱}
\begin{itemize}
    \item \textbf{除法运算}：
    \begin{itemize}
        \item \texttt{5 / 2} 结果为 \texttt{2}（整数除法截断小数）。
        \item \texttt{5.0 / 2} 结果为 \texttt{2.5}（浮点数参与运算）。
\begin{lstlisting}
int a = 5 / 2;   // a = 2
double b = 5.0 / 2; // b = 2.5
\end{lstlisting}
    \end{itemize}
    \item \textbf{赋值与相等}：
    \begin{itemize}
        \item \textbf{致命错误}：\texttt{if (a = 5)}。这会将 5 赋值给 a，且表达式结果为真（非0），导致条件永远成立。
        \item \textbf{正确写法}：\texttt{if (a == 5)}。
\begin{lstlisting}
if (a = 5) { ... } // 错误！永远为真，a 被改为 5
if (a == 5) { ... } // 正确
\end{lstlisting}
    \end{itemize}
    \item \textbf{浮点数比较}：
    \begin{itemize}
        \item \textbf{禁止}：\texttt{if (f == 0.0)}（因精度问题可能不准确）。
        \item \textbf{正确}：\texttt{if (fabs(f) < 1e-6)}。
\begin{lstlisting}
if (f == 0.0) ... // 危险！
if (fabs(f) < 1e-6) ... // 安全
\end{lstlisting}
    \end{itemize}
    \item \textbf{逻辑短路}：
    \begin{itemize}
        \item \texttt{A || B}：若 A 为真，B \textbf{不执行}。
        \item \texttt{A \&\& B}：若 A 为假，B \textbf{不执行}。
        \item \textbf{真题}：\texttt{k = (n=1) || (m=2);} $\rightarrow$ m 仍为原值，不会被赋值为 2。
\begin{lstlisting}
int n = 1, m = 0, k;
k = (n = 1) || (m = 2);
// n=1, k=1, m=0 (m=2 未执行)
\end{lstlisting}
    \end{itemize}
    \item \textbf{自增自减}：
    \begin{itemize}
        \item \texttt{b = a++}：先用 a 的旧值赋给 b，a 再自增。
        \item \texttt{b = ++a}：a 先自增，再把新值赋给 b。
\begin{lstlisting}
b = a++; // b=a, a=a+1
b = ++a; // a=a+1, b=a
\end{lstlisting}
    \end{itemize}
    \item \textbf{宏定义陷阱}：
    \begin{itemize}
        \item 宏只是简单的\textbf{文本替换}，不自动加括号。
        \item \texttt{\#define S(x) x*x} $\rightarrow$ \texttt{S(a+b)} 变为 \texttt{a+b*a+b}（错误）。
        \item \textbf{正确}：\texttt{\#define S(x) ((x)*(x))}。
\begin{lstlisting}
#define S(x) x*x
S(a+b) -> a+b*a+b // 错误
#define S(x) ((x)*(x))
S(a+b) -> ((a+b)*(a+b)) // 正确
\end{lstlisting}
    \end{itemize}
\end{itemize}

\subsection{数组与指针陷阱}
\begin{itemize}
    \item \textbf{字符串长度}：
    \begin{itemize}
        \item \texttt{sizeof("abc")} = 4 （包含末尾的 \texttt{'\textbackslash 0'}）。
        \item \texttt{strlen("abc")} = 3 （不包含 \texttt{'\textbackslash 0'}）。
\begin{lstlisting}
char s[] = "abc";
sizeof(s); // 4 ('\0')
strlen(s); // 3
\end{lstlisting}
    \end{itemize}
    \item \textbf{数组越界}：
    \begin{itemize}
        \item 定义 \texttt{int a[10];}，有效下标是 \texttt{0} 到 \texttt{9}。访问 \texttt{a[10]} 是越界！
\begin{lstlisting}
int a[10];
a[10] = 5; // 越界！最大下标是 9
\end{lstlisting}
    \end{itemize}
    \item \textbf{指针类型辨析}：
    \begin{itemize}
        \item \textbf{const 修饰}（口诀：左定值，右定向）：
        \begin{itemize}
            \item \texttt{const int *p}：\texttt{*p} (内容) 不可变，\texttt{p} (指向) 可变。
            \item \texttt{int * const p}：\texttt{p} (指向) 不可变，\texttt{*p} (内容) 可变。
        \end{itemize}
        \item \textbf{指针数组 vs 数组指针}：
        \begin{itemize}
            \item \texttt{int *p[10]}：\textbf{数组}，存了 10 个指针。
            \item \texttt{int (*p)[10]}：\textbf{指针}，指向长度为 10 的数组。
        \end{itemize}
        \item \textbf{函数指针}：
        \begin{itemize}
            \item 定义：\texttt{int (*func)(int, int);}
            \item 用途：回调函数（如 \texttt{qsort} 的比较函数）。
        \end{itemize}
        \item 详见附录 \hyperref[sec:appendix_ptr_types]{复杂指针类型代码详解}。
    \end{itemize}
    \item \textbf{野指针}：
    \begin{itemize}
        \item 定义指针时未初始化：\texttt{int *p;} $\rightarrow$ p 指向随机地址，直接 \texttt{*p = 10;} 会崩溃。
        \item \textbf{正确}：\texttt{int *p = NULL;} 或 \texttt{int a; int *p = \&a;}。
\begin{lstlisting}
int *p; *p = 10; // 崩溃！p 未初始化
int *p = NULL; // 安全
\end{lstlisting}
    \end{itemize}
\end{itemize}

\subsection{控制结构与输入输出陷阱}
\begin{itemize}
    \item \textbf{Switch 语句四大雷区 (真题解析)}：
    \begin{itemize}
        \item \textbf{雷区1：Case 标签必须是常量}
        \begin{itemize}
            \item 错误：\texttt{case c1:} (变量)。
            \item 正确：\texttt{case 1:} 或 \texttt{case 'A':} (整型/字符型常量表达式)。
        \end{itemize}
        \item \textbf{雷区2：Case 标签必须是整型}
        \begin{itemize}
            \item 错误：\texttt{case 1.0f:} (浮点数)。
            \item 正确：\texttt{switch} 表达式和 \texttt{case} 标签都必须是整型 (int, char, enum)。
        \end{itemize}
        \item \textbf{雷区3：语法结构错误}
        \begin{itemize}
            \item 错误：\texttt{switch(a);} (多了分号，导致 switch 无效)。
            \item 错误：\texttt{switch a} (少了括号)。
        \end{itemize}
        \item \textbf{雷区4：穿透性 (Fallthrough)}
        \begin{itemize}
            \item \texttt{case} 语句块结束后若没有 \texttt{break}，会直接执行下一个 \texttt{case} 的内容。
        \end{itemize}
    \end{itemize}
    \begin{lstlisting}
// 典型真题错误示例：
switch (a) {
    case c1: ... // 错！变量
    case 1.5: ... // 错！浮点
}
switch (a); { ... } // 错！分号截断
    \end{lstlisting}

    \item \textbf{悬挂 Else}：
    \begin{itemize}
        \item \texttt{else} 总是与\textbf{最近的未匹配的} \texttt{if} 配对，与缩进无关。
\begin{lstlisting}
if (a)
    if (b) x++;
else y++; // 这个 else 属于 if(b)
\end{lstlisting}
    \end{itemize}
    \item \textbf{分号错误}：
    \begin{itemize}
        \item \texttt{if (x > 0);} $\rightarrow$ 分号表示空语句，if 实际上什么也没控制。
\begin{lstlisting}
if (x > 0); // 多了分号
    x++;    // 无论 x 是否 > 0 都会执行
\end{lstlisting}
    \end{itemize}
    \item \textbf{scanf 缓冲区残留}：
    \begin{itemize}
        \item 先输入数字再输入字符时，回车符会被 \texttt{\%c} 读取。
        \item \textbf{解决}：\texttt{scanf(" \%c", \&ch);} （\%c 前加空格吃掉空白符）。
\begin{lstlisting}
scanf("%d", &n);
scanf(" %c", &ch); // 空格吃掉回车
\end{lstlisting}
    \end{itemize}
\end{itemize}

\subsection{函数参数传递陷阱}
\begin{itemize}
    \item \textbf{修改实参需传指针}：
    \begin{itemize}
        \item \textbf{原理}：C语言函数参数默认是\textbf{值传递} (Pass by Value)。形参只是实参的拷贝，修改形参不会影响实参。
        \item \textbf{正确做法}：传地址 (\texttt{\&var})，用指针接收 (\texttt{*ptr})，解引用修改 (\texttt{*ptr = val})。
        \item \textbf{常见错误}：直接传值，在函数内修改形参，以为实参也会变。
\begin{lstlisting}
// 错误：值传递，a, b 没变
void swap_wrong(int x, int y) { int t=x; x=y; y=t; }
// 正确：地址传递
void swap_right(int *x, int *y) { int t=*x; *x=*y; *y=t; }
// 调用：swap_right(&a, &b);
\end{lstlisting}
    \end{itemize}

    \item \textbf{数组传参退化为指针}：
    \begin{itemize}
        \item \textbf{原理}：数组名作为函数参数传递时，会\textbf{退化}为指向首元素的指针。
        \item \textbf{后果}：在函数内部无法通过 \texttt{sizeof(arr)} 获取数组总大小（只能得到指针大小，4或8字节）。
        \item \textbf{正确做法}：必须额外传递数组长度 \texttt{n}。
\begin{lstlisting}
void func(int arr[]) {
    // sizeof(arr) 在这里等于 sizeof(int*)，即 4 或 8
    // 无法知道数组真实长度！
}
// 正确：void func(int arr[], int n)
\end{lstlisting}
    \end{itemize}
\end{itemize}

\section{简答题预测汇总}

以下是根据历年考点和课程核心内容整理的简答题题库，建议背诵关键词。

\subsection{计算机基础与算法}
\begin{description}
    \item[Q1: 简述冯·诺依曼体系结构的核心思想（“存储程序”原理）。] \hfill \\
    \textbf{A:} 将程序（指令序列）和数据一起存储在计算机的存储器中，计算机在运行时自动地从存储器中取出指令并执行，而不需要人工干预。
    
    \item[Q2: 什么是算法？它有哪些基本特征？] \hfill \\
    \textbf{A:} 算法是解决特定问题的一系列明确、有限的指令。
    \textbf{五大特征}：
    \begin{enumerate}
        \item \textbf{有穷性}：步骤有限，在有限时间内完成。
        \item \textbf{确定性}：每一步指令清晰无歧义。
        \item \textbf{可行性}：每一步都能通过基本运算实现。
        \item \textbf{输入}：0个或多个输入。
        \item \textbf{输出}：1个或多个输出。
    \end{enumerate}

    \item[Q3: 结构化程序设计的三种基本控制结构是什么？] \hfill \\
    \textbf{A:} 顺序结构、选择结构（分支结构）、循环结构。
\end{description}

\subsection{C语言核心概念}
\begin{description}
    \item[Q4: 什么是指针？为什么要使用指针？] \hfill \\
    \textbf{A:} 指针是存储内存地址的变量。
    \textbf{作用}：
    \begin{enumerate}
        \item 可以直接操作内存地址，提高程序效率。
        \item 实现函数间的大数据传递（避免拷贝整个数组或结构体）。
        \item 支持动态内存分配（malloc/free）。
        \item 处理复杂的数据结构（如链表、树）。
    \end{enumerate}

    \item[Q5: 简述递归函数的优缺点。] \hfill \\
    \textbf{A:} 
    \textbf{优点}：代码简洁，逻辑清晰，适合解决具有重复子结构的问题（如汉诺塔、树的遍历）。
    \textbf{缺点}：函数调用开销大，占用栈空间多，深度过深可能导致栈溢出，效率通常低于迭代。

    \item[Q6: 结构体 (struct) 和共用体 (union) 的区别是什么？] \hfill \\
    \textbf{A:} 
    \textbf{结构体}：每个成员拥有独立的内存空间，总大小大于等于所有成员大小之和（受内存对齐影响）。所有成员可以同时存在。
    \textbf{共用体}：所有成员共享同一块内存空间，总大小等于最大成员的大小。同一时刻只能存储其中一个成员的值。

    \item[Q7: 局部变量和全局变量的区别？] \hfill \\
    \textbf{A:} 
    \textbf{作用域}：局部变量仅在定义它的函数或代码块内有效；全局变量在整个源程序文件中有效。
    \textbf{生命周期}：局部变量随函数调用创建，函数结束销毁（static除外）；全局变量在程序运行期间一直存在。
    \textbf{优先级}：同名时，局部变量屏蔽全局变量。

    \item[Q8: 关键字 static 的作用有哪些？] \hfill \\
    \textbf{A:} 
    \begin{enumerate}
        \item \textbf{修饰局部变量}：改变生命周期，使其在程序运行期间一直存在，但作用域不变（仍局限于函数内）。
        \item \textbf{修饰全局变量/函数}：限制作用域，使其仅在当前源文件中可见，不能被其他文件引用（隐藏）。
    \end{enumerate}
\end{description}

\subsection{常见算法特点与比较}
\begin{description}
    \item[Q9: 简述三种基本排序算法（冒泡、选择、插入）的区别与特点。] \hfill \\
    \textbf{A:} 
    \begin{itemize}
        \item \textbf{冒泡排序}：两两比较相邻元素，大数下沉。
        \begin{itemize}
            \item 最好 $O(n)$（已有序），最坏 $O(n^2)$。
            \item \textbf{稳定}。
        \end{itemize}
        \item \textbf{选择排序}：每次从未排序区选最小的放到已排序区末尾。
        \begin{itemize}
            \item 无论是否有序，复杂度恒为 $O(n^2)$。
            \item \textbf{不稳定}（如序列 5, 8, 5, 2，第一趟交换第一个5和2，导致两个5相对位置改变）。
        \end{itemize}
        \item \textbf{插入排序}：将新元素插入到已排序序列的合适位置。
        \begin{itemize}
            \item 最好 $O(n)$，最坏 $O(n^2)$。对于\textbf{基本有序}的数据效率最高。
            \item \textbf{稳定}。
        \end{itemize}
    \end{itemize}

    \item[Q10: 顺序查找与二分查找的对比。] \hfill \\
    \textbf{A:} 
    \begin{itemize}
        \item \textbf{顺序查找}：
        \begin{itemize}
            \item 适用性：无序或有序数组均可，链表也可。
            \item 复杂度：$O(n)$。
        \end{itemize}
        \item \textbf{二分查找}：
        \begin{itemize}
            \item 适用性：\textbf{必须是有序数组}（支持随机访问）。
            \item 复杂度：$O(\log_2 n)$。效率远高于顺序查找。
        \end{itemize}
    \end{itemize}

    \item[Q11: 什么是算法的时间复杂度和空间复杂度？] \hfill \\
    \textbf{A:} 
    \begin{itemize}
        \item \textbf{时间复杂度}：算法执行所需时间与问题规模 $n$ 之间的增长关系（通常用大O表示法，如 $O(n), O(n^2)$）。衡量算法运行快慢。
        \item \textbf{空间复杂度}：算法执行过程中所需的存储空间与问题规模 $n$ 之间的增长关系。衡量算法占用内存多少。
    \end{itemize}
\end{description}

\subsection{实验课补充简答题}
\begin{description}
    \item[Q12: 有序数组与无序数组在操作效率上有何区别？] \hfill \\
    \textbf{A:}
    \begin{itemize}
        \item \textbf{无序数组}：插入最快 $O(1)$（直接追加），查找和删除较慢 $O(n)$（线性查找）。
        \item \textbf{有序数组}：查找最快 $O(\log n)$（二分查找），插入和删除较慢 $O(n)$（需要移动元素保持有序）。
    \end{itemize}

    \item[Q13: 栈 (Stack) 的主要特性和基本操作有哪些？] \hfill \\
    \textbf{A:}
    \begin{itemize}
        \item \textbf{特性}：后进先出 (LIFO, Last In First Out)。
        \item \textbf{基本操作}：
        \begin{itemize}
            \item \textbf{Push (入栈)}：将元素放入栈顶。
            \item \textbf{Pop (出栈)}：移除并返回栈顶元素。
            \item \textbf{Peek (查看栈顶)}：返回栈顶元素但不移除。
            \item \textbf{IsEmpty (判空)}：检查栈是否为空。
        \end{itemize}
    \end{itemize}

    \item[Q14: 如何利用栈实现括号匹配（如JSON字符串校验）？] \hfill \\
    \textbf{A:} 遍历字符串，遇到左括号（如 \texttt{\{, [}）入栈；遇到右括号（如 \texttt{\}, ]}）时，检查栈顶元素是否为对应的左括号。若是，则弹出栈顶元素继续；若否或栈为空，则匹配失败。遍历结束后栈应为空。

    \item[Q15: 在“找出兴趣值最大的k个客户”问题中，为什么不需要全排序？] \hfill \\
    \textbf{A:} 全排序的时间复杂度通常为 $O(n \log n)$ 或 $O(n^2)$。如果只需前 $k$ 个最大值，可以使用\textbf{部分选择排序}（执行 $k$ 轮选择），时间复杂度为 $O(k \cdot n)$。当 $k \ll n$ 时，效率更高。

    \item[Q16: 递归解决汉诺塔问题的核心逻辑是什么？] \hfill \\
    \textbf{A:} 将 $n$ 个盘子从 A 移到 C 的过程分解为三步：
    \begin{enumerate}
        \item 将 $n-1$ 个盘子从 A 移到 B（借助 C）。
        \item 将第 $n$ 个盘子（最大的）从 A 移到 C。
        \item 将 $n-1$ 个盘子从 B 移到 C（借助 A）。
    \end{enumerate}
    这是一个典型的分治递归思想。

    \item[Q17: 计算机五大部件（运算器、控制器、存储器、输入设备、输出设备）是什么？] \hfill \\
    \textbf{A:} 这五大部件构成了冯·诺依曼体系结构的基础。
    \begin{itemize}
        \item \textbf{运算器}：负责算术和逻辑运算。
        \item \textbf{控制器}：指挥各部件协调工作。
        \item \textbf{存储器}：存放程序和数据。
        \item \textbf{输入/输出设备}：负责与外部交互。
    \end{itemize}
    详细的原理介绍及结构图请跳转至 \hyperref[sec:von_neumann]{附录：冯·诺依曼体系结构详解}。
\end{description}

\appendix
\section{冯·诺依曼体系结构详解} \label{sec:von_neumann}

计算机硬件系统主要由五大部件组成：运算器、控制器、存储器、输入设备和输出设备。这种结构被称为\textbf{冯·诺依曼体系结构}。

\begin{enumerate}
    \item \textbf{运算器 (Arithmetic Logic Unit, ALU)}: 
    \begin{itemize}
        \item \textbf{原理}：计算机的“大脑”之一，负责执行所有的算术运算（如加、减、乘、除）和逻辑运算（如与、或、非、异或）。
        \item \textbf{作用}：对数据进行加工处理。
    \end{itemize}
    
    \item \textbf{控制器 (Control Unit, CU)}: 
    \begin{itemize}
        \item \textbf{原理}：计算机的“指挥中心”。它从存储器中取出指令，进行译码分析，然后向其他部件发出控制信号。
        \item \textbf{作用}：保证计算机各部件有序、协调地工作。通常与运算器集成在一起，统称为\textbf{中央处理器 (CPU)}。
    \end{itemize}

    \item \textbf{存储器 (Memory)}: 
    \begin{itemize}
        \item \textbf{原理}：具有记忆功能的部件，由许多存储单元组成，每个单元都有唯一的地址。
        \item \textbf{作用}：用于存放程序（指令序列）和数据。分为\textbf{内存储器}（速度快，CPU直接访问）和\textbf{外存储器}（容量大，长期保存）。
    \end{itemize}

    \item \textbf{输入设备 (Input Device)}: 
    \begin{itemize}
        \item \textbf{原理}：将外部信息（文字、声音、图像等）转换为计算机能识别的二进制数据。
        \item \textbf{作用}：向计算机输入原始数据和程序。例如：键盘、鼠标、扫描仪。
    \end{itemize}

    \item \textbf{输出设备 (Output Device)}: 
    \begin{itemize}
        \item \textbf{原理}：将计算机处理后的二进制结果转换为人类能识别的形式。
        \item \textbf{作用}：输出处理结果。例如：显示器、打印机、音响。
    \end{itemize}
\end{enumerate}

\subsection*{结构示意图}

\begin{center}
\begin{tikzpicture}[node distance=1.5cm]
    % Styles
    \tikzstyle{block} = [rectangle, draw, fill=blue!10, text width=5em, text centered, rounded corners, minimum height=3em]
    \tikzstyle{bigblock} = [rectangle, draw, fill=red!10, text width=8em, text centered, rounded corners, minimum height=6em]
    \tikzstyle{line} = [draw, -latex', thick]

    % Nodes
    \node [block] (input) {输入设备};
    \node [bigblock, right=1.5cm of input] (cpu) {CPU\\(运算器+控制器)};
    \node [block, right=1.5cm of cpu] (output) {输出设备};
    \node [block, above=1.5cm of cpu] (memory) {存储器};

    % Paths
    \path [line] (input) -- (cpu);
    \path [line] (cpu) -- (output);
    \path [line] (cpu) -- (memory);
    \path [line] (memory) -- (cpu);
    
    % Annotations
    \node [below=0.2cm of input] {\footnotesize 数据输入};
    \node [below=0.2cm of output] {\footnotesize 结果输出};
\end{tikzpicture}
\end{center}

\newpage
\section{附录 1：程序填空题常考模式代码详解}
\label{sec:appendix_code}

\subsection{数学公式实现}
\begin{itemize}
    \item \textbf{海伦公式求三角形面积}
    \begin{lstlisting}
#include <math.h>
// ...
double a, b, c, s, area;
// 输入 a, b, c
s = (a + b + c) / 2.0;
area = sqrt(s * (s - a) * (s - b) * (s - c));
    \end{lstlisting}
    
    \item \textbf{泰勒级数求 sin(x)}
    \begin{lstlisting}
#include <math.h>
// ...
double x, term, sum;
int n = 1;
// 输入 x
term = x;
sum = x;
// 循环终止条件：某一项绝对值小于 1e-5
while (fabs(term) >= 1e-5) {
    term = -term * x * x / ((n + 1) * (n + 2));
    sum += term;
    n += 2;
}
    \end{lstlisting}
\end{itemize}

\subsection{最大公约数 (辗转相除法)}
\begin{lstlisting}
int a, b, temp;
// 输入 a, b
while (b != 0) {
    temp = a % b;
    a = b;
    b = temp;
}
// 循环结束后，a 即为最大公约数
\end{lstlisting}

\subsection{条件判断 (三目运算符嵌套)}
\begin{lstlisting}
// 求 a, b, c 中的最大值
int max = (a > b) ? ((a > c) ? a : c) : ((b > c) ? b : c);
\end{lstlisting}

\section{附录 2：标准算法模板}
\label{sec:appendix_algo}

\subsection{排序算法}
\begin{itemize}
    \item \textbf{冒泡排序 (Bubble Sort)}
    \begin{lstlisting}
void bubble_sort(int arr[], int n) {
    int i, j, temp;
    for (i = 0; i < n - 1; i++) {
        for (j = 0; j < n - 1 - i; j++) {
            if (arr[j] > arr[j + 1]) {
                temp = arr[j];
                arr[j] = arr[j + 1];
                arr[j + 1] = temp;
            }
        }
    }
}
    \end{lstlisting}
    \item \textbf{选择排序 (Selection Sort)}
    \begin{lstlisting}
void selection_sort(int arr[], int n) {
    int i, j, min_idx, temp;
    for (i = 0; i < n - 1; i++) {
        min_idx = i;
        for (j = i + 1; j < n; j++) {
            if (arr[j] < arr[min_idx]) min_idx = j;
        }
        if (min_idx != i) {
            temp = arr[i];
            arr[i] = arr[min_idx];
            arr[min_idx] = temp;
        }
    }
}
    \end{lstlisting}
\end{itemize}

\subsection{查找算法}
\begin{itemize}
    \item \textbf{二分查找 (Binary Search)}
    \begin{lstlisting}
int binary_search(int arr[], int n, int key) {
    int low = 0, high = n - 1, mid;
    while (low <= high) {
        mid = low + (high - low) / 2;
        if (arr[mid] == key) return mid;
        else if (arr[mid] < key) low = mid + 1;
        else high = mid - 1;
    }
    return -1;
}
    \end{lstlisting}
\end{itemize}

\subsection{其他常用算法}
\begin{itemize}
    \item \textbf{素数判断}
    \begin{lstlisting}
int is_prime(int n) {
    if (n <= 1) return 0;
    int limit = (int)sqrt(n);
    for (int i = 2; i <= limit; i++) {
        if (n % i == 0) return 0;
    }
    return 1;
}
    \end{lstlisting}
    \item \textbf{字符串逆序}
    \begin{lstlisting}
void reverse_string(char str[]) {
    int len = strlen(str);
    int i = 0, j = len - 1;
    char temp;
    while (i < j) {
        temp = str[i]; str[i] = str[j]; str[j] = temp;
        i++; j--;
    }
}
    \end{lstlisting}
\end{itemize}

\section{附录 3：复杂指针类型深度解析}
\label{sec:appendix_ptr_types}

\subsection{1. const 与指针的爱恨情仇}
\textbf{核心原理}：\texttt{const} 修饰的是紧跟在它后面的那个“东西”。
\begin{itemize}
    \item \textbf{常量指针 (Pointer to Constant)}：\texttt{const int *p}
    \begin{itemize}
        \item \textbf{解析}：\texttt{const} 修饰的是 \texttt{int} (内容)。
        \item \textbf{理解}：\texttt{*p} 被锁死了，不能通过 \texttt{p} 修改它指向的数据。但 \texttt{p} 自己没锁，可以改指别人。
        \item \textbf{场景}：函数参数 \texttt{void func(const char *str)}，保证函数内部不会修改传入的字符串。
    \end{itemize}
    \item \textbf{指针常量 (Constant Pointer)}：\texttt{int * const p}
    \begin{itemize}
        \item \textbf{解析}：\texttt{const} 修饰的是 \texttt{p} (指针变量本身)。
        \item \textbf{理解}：\texttt{p} 被锁死了，一旦初始化就不能指向别处。但 \texttt{*p} (指向的内容) 是自由的。
        \item \textbf{场景}：硬件驱动中固定地址的寄存器指针。
    \end{itemize}
\end{itemize}
\begin{lstlisting}
int a = 10, b = 20;
const int *p1 = &a; // "我只读，不改内容"
// *p1 = 30; // Error: read-only location
p1 = &b;     // OK: p1 可以变心

int * const p2 = &a; // "我专一，只指你"
*p2 = 30;    // OK: 内容可以变
// p2 = &b;  // Error: assignment of read-only variable
\end{lstlisting}

\textbf{核心区别对比表}：
\begin{center}
\begin{tabular}{|c|c|c|}
\hline
\textbf{特性} & \textbf{const int *p} & \textbf{int * const p} \\
\hline
名称 & 常量指针 & 指针常量 \\
\hline
被锁定的对象 & 内容 (\texttt{*p}) & 指针本身 (\texttt{p}) \\
\hline
可变的对象 & 指针本身 (\texttt{p}) & 内容 (\texttt{*p}) \\
\hline
口诀 & 左定值 & 右定向 \\
\hline
\end{tabular}
\end{center}

\subsection{2. 指针数组 vs 数组指针}
\textbf{核心原理}：优先级 \texttt{[]} > \texttt{*}。
\begin{itemize}
    \item \textbf{指针数组}：\texttt{int *p[10]}
    \begin{itemize}
        \item \textbf{解析}：\texttt{p} 先和 \texttt{[10]} 结合 $\rightarrow$ \texttt{p} 是个数组。数组里存什么？存 \texttt{int *}。
        \item \textbf{内存模型}：占 $10 \times 4/8$ 字节（取决于指针大小）。
        \item \textbf{场景}：字符串数组 \texttt{char *names[] = \{"Tom", "Jerry"\};}，处理不规则长度的字符串列表。
    \end{itemize}
    \item \textbf{数组指针}：\texttt{int (*p)[10]}
    \begin{itemize}
        \item \textbf{解析}：\texttt{(*p)} 强制 \texttt{p} 先是因指针。指向什么？指向 \texttt{int[10]} 这种类型的数组。
        \item \textbf{内存模型}：占 4/8 字节（只是一个指针）。
        \item \textbf{场景}：处理二维数组。二维数组名 \texttt{arr} 传参时退化为数组指针。
    \end{itemize}
\end{itemize}
\begin{lstlisting}
// 指针数组：本质是"数组"
char *strs[3] = {"Apple", "Banana", "Cherry"};
// strs[0] 是一个指针，指向 "Apple" 的首地址

// 数组指针：本质是"指针"
int matrix[2][3] = {{1,2,3}, {4,5,6}};
int (*p)[3] = matrix; // p 指向 matrix 的第 0 行 (包含3个int的数组)
// 访问 matrix[1][2] 的写法：
// 1. p[1][2]
// 2. *(*(p+1)+2)
\end{lstlisting}

\textbf{核心区别对比表}：
\begin{center}
\begin{tabular}{|c|c|c|}
\hline
\textbf{特性} & \textbf{int *p[10]} & \textbf{int (*p)[10]} \\
\hline
名称 & 指针数组 & 数组指针 \\
\hline
本质 & 数组 & 指针 \\
\hline
存储内容 & 10 个 \texttt{int*} 指针 & 1 个指向数组的地址 \\
\hline
内存占用 & $10 \times$ 指针大小 & 1 $\times$ 指针大小 \\
\hline
典型用途 & 字符串数组 & 二维数组传参 \\
\hline
\end{tabular}
\end{center}

\subsection{3. 函数指针与回调机制}
\textbf{核心原理}：函数名在表达式中被解读为该函数的入口地址。
\begin{itemize}
    \item \textbf{定义}：\texttt{int (*func)(int, int)}
    \begin{itemize}
        \item \texttt{(*func)}：表明 \texttt{func} 是个指针。
        \item \texttt{(int, int)}：表明它指向的函数接受两个 int 参数。
        \item \texttt{int}：表明它指向的函数返回 int。
    \end{itemize}
    \item \textbf{qsort 深度解析}：
    \begin{itemize}
        \item \texttt{void *}：通用指针，可以接收任何类型的地址，但不能直接解引用。
        \item \textbf{类型转换}：\texttt{*(int*)a}。先将 \texttt{void*} 强转为 \texttt{int*} (告诉编译器这是个整数地址)，再 \texttt{*} 取值。
    \end{itemize}
\end{itemize}
\begin{lstlisting}
// 比较函数：返回 <0, 0, >0
int compare(const void *a, const void *b) {
    // 1. (int*)a : 把 void指针 当作 int指针 看待
    // 2. *(int*)a : 取出该地址里的整数值
    return (*(int*)a - *(int*)b); 
}

int main() {
    int values[] = { 88, 56, 100, 2, 25 };
    // qsort 内部会多次调用 compare 函数来决定顺序
    // 这就是"回调"：你把逻辑传给 qsort，qsort 回头调用你的逻辑
    qsort(values, 5, sizeof(int), compare);
}
\end{lstlisting}

\subsection{4. 二维数组 vs 二维动态数组 (内存模型与访问)}
\textbf{核心原理}：内存连续性决定了访问方式的本质不同。

\begin{itemize}
    \item \textbf{静态二维数组}：\texttt{int arr[M][N]}
    \begin{itemize}
        \item \textbf{内存布局}：\textbf{线性连续}。所有元素像糖葫芦一样串在一起。
        \item \textbf{寻址公式}：\texttt{Addr = Base + (i * N + j) * sizeof(int)}
        \item \textbf{过程}：
        \begin{enumerate}
            \item CPU 拿到基地址 \texttt{Base}。
            \item 编译器直接生成计算指令：$i \times N + j$。
            \item 一次性算出最终地址，直接访问。
        \end{enumerate}
        \item \textbf{特点}：速度快，但列数 $N$ 必须在编译时确定（或作为类型信息存在）。
    \end{itemize}
    
    \item \textbf{动态二维数组 (二级指针)}：\texttt{int **p}
    \begin{itemize}
        \item \textbf{内存布局}：\textbf{离散破碎}。有一个“脊柱”数组存指针，指向多根“肋骨”数组存数据。
        \item \textbf{寻址公式}：\texttt{Addr = *(*p + i) + j * sizeof(int)} (伪代码)
        \item \textbf{过程}：
        \begin{enumerate}
            \item \textbf{第一次访存}：访问 \texttt{p[i]} (即 \texttt{*(p+i)})，从内存中读取第 $i$ 行的首地址 \texttt{RowAddr}。
            \item \textbf{第二次访存}：在 \texttt{RowAddr} 的基础上偏移 $j$，访问最终数据。
        \end{enumerate}
        \item \textbf{特点}：灵活（每行长度可以不同），但多一次内存访问，且缓存命中率 (Cache Hit) 可能较低。
    \end{itemize}
\end{itemize}

\begin{lstlisting}
// 1. 静态二维数组
int arr[3][4]; 
// arr[1][2] -> 编译器直接算出偏移量：1*4 + 2 = 6 -> 访问第6个元素

// 2. 动态二维数组 (二级指针)
int **p = (int**)malloc(3 * sizeof(int*)); // 脊柱 (存行指针)
for(int i=0; i<3; i++) {
    p[i] = (int*)malloc(4 * sizeof(int));  // 肋骨 (存数据)
}
// p[1][2] -> CPU先去 p[1] 查到地址 0x5000，再去 0x5000 + 2*4 取值
\end{lstlisting}

\textbf{核心区别对比表}：
\begin{center}
\begin{tabular}{|c|c|c|}
\hline
\textbf{特性} & \textbf{int arr[M][N]} & \textbf{int **p} \\
\hline
内存布局 & 连续的一整块 & 破碎的 (指针数组+多块数据) \\
\hline
寻址方式 & 一次计算 (基址+偏移) & 两次访问 (查表+偏移) \\
\hline
类型 & \texttt{int (*)[N]} (数组指针) & \texttt{int **} (二级指针) \\
\hline
传参写法 & \texttt{func(int a[][N])} & \texttt{func(int **a)} \\
\hline
\end{tabular}
\end{center}

\subsection{5. 真题解析：二维数组指针表示法 (a[i][j] 的 N 种写法)}
\textbf{题目}：定义二维数组 \texttt{int a[m][n]}，以下表示 \texttt{a[i][j]} 的方式正确的是？
\begin{itemize}
    \item A. \texttt{*(a + i*n + j)}
    \item C. \texttt{*(*(a + i) + j)}
\end{itemize}

\textbf{深度辨析}：
\begin{itemize}
    \item \textbf{选项 C (正确)}：\texttt{*(*(a + i) + j)}
    \begin{enumerate}
        \item \texttt{a}：类型是 \texttt{int (*)[n]} (行指针)。
        \item \texttt{a + i}：指向第 $i$ 行 (地址步长为 $n \times 4$)。
        \item \texttt{*(a + i)}：解引用行指针，得到第 $i$ 行 (即一维数组名 \texttt{a[i]})。此时类型退化为 \texttt{int *} (列指针)。
        \item \texttt{*(a + i) + j}：在第 $i$ 行内偏移 $j$ 个元素。
        \item \texttt{*(\dots)}：取出最终的整数值。
    \end{enumerate}
    
    \item \textbf{选项 A \& B (错误原因深度剖析)}：
    \begin{itemize}
        \item \textbf{选项 A}：\texttt{*(a + i*n + j)}
        \begin{itemize}
            \item \textbf{错误本质}：\textbf{步长错误}。
            \item \texttt{a} 是行指针，\texttt{a+1} 的含义是“下一行”。
            \item \texttt{a + (i*n + j)} 意味着跳过了 $i \times n + j$ \textbf{行}！这会访问到极其遥远的非法内存区域。
            \item \textbf{修正}：必须强转为 \texttt{int*} 才能按“个”跳。即 \texttt{*((int*)a + i*n + j)}。
        \end{itemize}
        \item \textbf{选项 B}：\texttt{*(a + j*n + i)}
        \begin{itemize}
            \item \textbf{错误本质}：\textbf{步长错误 + 逻辑混乱}。
            \item 同样未强转，导致按“行”跳跃。
            \item 即使强转了，公式 $j \times n + i$ 也是错误的（这是列优先存储的计算方式，而 C 语言是行优先存储）。
        \end{itemize}
    \end{itemize}
\end{itemize}

\textbf{扩展：静态数组 vs 动态数组的写法通用性}
\begin{center}
\begin{tabular}{|c|c|c|}
\hline
\textbf{写法} & \textbf{静态数组 \texttt{int a[M][N]}} & \textbf{动态数组 \texttt{int **p}} \\
\hline
\texttt{a[i][j]} & \textcolor{green}{\checkmark} & \textcolor{green}{\checkmark} \\
\hline
\texttt{*(*(a+i)+j)} & \textcolor{green}{\checkmark} (行指针 $\to$ 列指针) & \textcolor{green}{\checkmark} (二级指针 $\to$ 一级指针) \\
\hline
\texttt{*(a[i]+j)} & \textcolor{green}{\checkmark} & \textcolor{green}{\checkmark} \\
\hline
\texttt{(*(a+i))[j]} & \textcolor{green}{\checkmark} & \textcolor{green}{\checkmark} \\
\hline
\texttt{*(\&a[0][0] + i*N + j)} & \textcolor{green}{\checkmark} (内存连续) & \textcolor{red}{\ding{55}} (内存不连续) \\
\hline
\end{tabular}
\end{center}
\textbf{结论}：前四种写法是通用的（语法糖不同，底层逻辑不同但表现一致）。唯独\textbf{线性计算法}（最后一行）只适用于内存连续的静态数组。

\section*{附录4：字符与字符串陷阱深度解析}
\label{sec:appendix_char_traps}

\subsection{1. 转义字符 (Escape Characters) 全解}
\textbf{核心定义}：以反斜杠 \texttt{\textbackslash} 开头的特殊字符序列。

\begin{itemize}
    \item \textbf{常见转义字符表}：
    \begin{center}
    \begin{tabular}{|c|l|c|l|}
    \hline
    \textbf{转义字符} & \textbf{含义} & \textbf{ASCII} & \textbf{备注} \\
    \hline
    \texttt{\textbackslash n} & 换行 (Newline) & 10 & 光标移到下一行开头 \\
    \hline
    \texttt{\textbackslash r} & 回车 (Carriage Return) & 13 & 光标移到本行开头 \\
    \hline
    \texttt{\textbackslash t} & 水平制表符 (Tab) & 9 & 光标移到下一个制表位 \\
    \hline
    \texttt{\textbackslash b} & 退格 (Backspace) & 8 & 光标左移一格 \\
    \hline
    \texttt{\textbackslash 0} & 空字符 (Null) & 0 & 字符串结束标志 \\
    \hline
    \texttt{\textbackslash \textbackslash} & 反斜杠本身 & 92 & \\
    \hline
    \texttt{\textbackslash '} & 单引号 & 39 & 用于字符常量 \texttt{'\textbackslash ''} \\
    \hline
    \texttt{\textbackslash "} & 双引号 & 34 & 用于字符串常量 \texttt{"\textbackslash ""} \\
    \hline
    \end{tabular}
    \end{center}
\end{itemize}

\subsection{2. 数值转义：\textbackslash ddd vs \textbackslash xhh}
\textbf{核心区别对比表}：
\begin{center}
\begin{tabular}{|c|c|c|}
\hline
\textbf{特性} & \textbf{\textbackslash ddd (八进制)} & \textbf{\textbackslash xhh (十六进制)} \\
\hline
前缀 & \texttt{\textbackslash} 后直接跟数字 & \texttt{\textbackslash x} \\
\hline
有效字符 & 0-7 & 0-9, a-f, A-F \\
\hline
\textbf{长度限制} & \textbf{最多 3 位} & \textbf{无限制 (贪婪匹配)} \\
\hline
示例 & \texttt{'\textbackslash 101'} ('A') & \texttt{'\textbackslash x41'} ('A') \\
\hline
\end{tabular}
\end{center}

\textbf{高频陷阱解析}：
\begin{itemize}
    \item \textbf{陷阱1：非法八进制}
    \begin{itemize}
        \item \texttt{'\textbackslash 08'} 是非法的！因为 8 不是八进制数。
        \item 编译器会报错，或者将其解析为 \texttt{'\textbackslash 0'} 和 \texttt{'8'} (如果是字符串)。
    \end{itemize}
    
    \item \textbf{陷阱2：八进制截断}
    \begin{itemize}
        \item 题目：\texttt{strlen("\textbackslash 128")} 是多少？
        \item 解析：\texttt{\textbackslash 12} 是一个八进制转义 (ASCII 10)，\texttt{8} 是普通字符。
        \item 结果：长度为 2。
    \end{itemize}

    \item \textbf{陷阱3：十六进制贪婪}
    \begin{itemize}
        \item 题目：\texttt{strlen("\textbackslash x123")} 是多少？
        \item 解析：编译器会一直读下去，试图解析 \texttt{0x123} (291)，这超出了 \texttt{char} (0-255) 的范围。
        \item 结果：通常报错或发生截断。
        \item \textbf{对策}：如果要表示 \texttt{\textbackslash x12} 和字符 \texttt{3}，必须写成 \texttt{"\textbackslash x12" "3"} (利用字符串自动拼接)。
    \end{itemize}

    \item \textbf{陷阱4：\textbackslash x vs 0x}
    \begin{itemize}
        \item \texttt{char c = '\textbackslash x41';} (正确，字符 'A')
        \item \texttt{int n = 0x41;} (正确，整数 65)
        \item \texttt{char c = '\textbackslash 0x41';} (\textbf{错误}！这是 \texttt{'\textbackslash 0'} + \texttt{'x'} + ...)
    \end{itemize}
\end{itemize}

\end{document}
